\documentclass[10pt,twocolumn]{article}
\usepackage[letterpaper,margin=.72in,columnsep=.23in]{geometry}
\usepackage{amsmath,graphicx,hyperref}
\IfFileExists{paper/metrics.tex}{\newcommand{\paperdir}{paper/}}{\newcommand{\paperdir}{}}
\input{\paperdir metrics.tex}
\newcommand{\placeholder}[1]{\texttt{[#1]}}
\title{Test-Time Compute in SmolVLA: A100 System Characterization and Early-Window Selection}
\author{\placeholder{Author list}\\\placeholder{Affiliations}}
\date{\placeholder{Target venue and submission format}}
\renewcommand{\dbltopfraction}{.85}
\setcounter{dbltopnumber}{3}
\raggedbottom
\begin{document}
\maketitle

\begin{abstract}
Frozen VGAS four-proposal selection on SmolVLA improved observed paired success over a single proposal on LIBERO-Object and LIBERO-Long Task2. Using four proposals only at early queries retained observed full-selection success within a 3 percentage point tolerance on new Object starts and Long Task2, cutting four-way calls by 51.5\% and 58.6\%, respectively. Across 500 paired Object layouts, first-five outperformed Late (+7.60 pp [2.60, 12.60]), with 2,500 versus 2,893 realized four-way calls.
On LIBERO-Goal, the preregistered one-third rule met both point-estimate criteria.
Object calibration EWS averaged 2,201 ms of core inference per episode versus 2,282 ms for one proposal. On A100, prefix sharing reduced rollout mean decision time from 272 to 217 ms, CUDA Graph replay reduced a separate matched fixed-observation median from 226 to 81 ms, and per-channel INT8 packing reduced persistent critic storage from 175 to 44 MB.
\end{abstract}

\section{Introduction}
A vision-language-action policy turns camera images, an instruction and robot
state into a sequence of actions. Under limited adaptation data, different
proposals from the same actor can lead to different outcomes. A critic can
choose among them. An efficient implementation must account for both proposal
generation and scoring. The observation stays fixed within a query, while each
proposal starts from independent noise.

Using the released SmolVLA actor and VGAS critic \cite{smolvla,vgas}, we observe paired success gains from four-way selection on Object and Long Task2 in the same LIBERO simulator. Fixed early windows retain the full schedule's observed success within the prespecified point-estimate tolerance while cutting scoring calls by about half, and the Late control measures the effect of delaying selection on Object. Across 500 paired Object layouts, first-five outperformed Late (+7.60 pp [2.60, 12.60]), with 2,500 versus 2,893 realized four-way calls.
Our A100 measurements characterize shared-prefix inference and matched CUDA Graph replay, and our separate INT8 evaluation measures the reduction in persistent critic storage.

\section{Related Work}
Value-Guided Policy Steering reranks actions using an offline-trained value
function \cite{vgps}. RoboMonkey studies test-time sampling and verification
for robot policies \cite{robomonkey}. VGAS provides the action-chunk critic
and scoring rule that we use \cite{vgas}. Q-Guided Flow and Guided Action
Flow use critic gradients during action generation \cite{qgf,gaf}.
Our evaluation compares completed chunks under a frozen selection rule.

Zora shares attention keys and values for autoregressive proposal verification.
PACT-VLA shares context while scoring proposals under a newer observation
\cite{zora,pact}. Adaptive decoding in diffusion driving VLAs targets the
denoising budget itself \cite{early}. CUDA Graphs can reduce repeated host
dispatch by replaying captured diffusion workloads \cite{graphs}. These
precedents motivate our measurement of how execution scheduling and proposal
batch size interact when a flow actor feeds a learned critic.

ELASTIC learns how to vary denoising steps and candidate count with
robot state \cite{elastic}. EWS allocates a fixed-size candidate pool
across queries in an episode and can be combined with within-query
compute allocation.

\section{Method}
Figure~\ref{fig:overview} follows one control query from observation to execution.
\begin{figure*}[t]
\centering
\includegraphics[width=\textwidth]{\paperdir figures/method-overview.pdf}
\caption{Shared-prefix selection at inference. Camera views come from a released
LIBERO-Object query. One prefix supplies both the actor's attention cache and the
critic's observation features. Four independent noise draws produce four action
chunks. Scores are the lower of two critic estimates from the recorded query. The lower panel shows the horizons used by the released checkpoints and
our executor.}
\label{fig:overview}
\end{figure*}
\input{\paperdir algorithm.tex}

\input{\paperdir tables/closed-loop.tex}

\section{System Design and Measurement}
The shared encoding produces attention keys and values for the actor and
observation features for the twin critic. Each candidate retains its own
denoising trajectory through the Euler updates. The critic evaluates the pool
as a batch and the executor slices out and postprocesses the selected candidate.
A new observation refreshes the prefix.

For the fixed-shape experiment we capture the denoising loop and refresh its KV
cache, mask, noise and time scalars before each replay. The surrounding decision
remains eager. Full-decision timing includes encoding, generation, scoring and
selection. The fixed-observation replay scope excludes preprocessing,
postprocessing, simulation and one-time capture. The experiment alternates
matched eager and replay runs after warm-up. Strict FP32 checks disable TF32
and compare returned actions, critic scores and winner indices.

\section{Evaluation}
\subsection{Closed-loop outcomes}
We pair by physical layout and policy noise. Layout-level resampling keeps
repeated noises together within a physical layout. Table~\ref{tab:closed-loop}
compares shared Bo4 with the single-proposal and same-width original comparators.
The lighting condition is a secondary perturbation check. Its familywise
interval spans zero. The confirmed Object cohort comes from the corrected
executor. An earlier wrapper always dispatched candidate zero, so that batch
is excluded from the evidence.

\input{\paperdir tables/scaling.tex}

\subsection{Early-window scheduled selection}
\label{sec:scheduled}
We keep the actor, critic and execution horizon fixed. As an episode
progresses, fewer control steps remain to recover from an earlier
decision. This motivates spending the selection budget early. We call this rule early-window selection (EWS). At queries
$1,\ldots,K$ it samples four proposals and ranks them with VGAS.
Subsequent queries use one proposal and skip the critic. Each query
executes up to 20 environment steps. Episodes allow up to 14 queries
on Object and 26 on Long Task2.

\input{\paperdir tables/scheduled-selection.tex}
\begin{figure}[tb]
\centering
\includegraphics[width=\columnwidth]{\paperdir figures/scheduled-dose.pdf}
\caption{LIBERO-Object dose sweep on 300 paired layouts. Error bars are
descriptive paired 95\% intervals. The dashed line scales the full
schedule's gain by $K/14$. The random-seven policy samples seven
positions before acting. Late schedules nine positions starting at query 6. Both stop on termination.}
\label{fig:scheduled-dose}
\end{figure}

Table~\ref{tab:scheduled} summarizes the paired outcomes. On the
300-layout calibration cohort, we chose the smallest tested
$K$ whose success-rate point estimate fell within 3 percentage points
of full selection. This is EWS's cutoff calibration procedure. This cohort overlaps with a 50-layout pilot that
motivated the sweep. Gains flattened after the first few queries
(Fig.~\ref{fig:scheduled-dose}). The first five reached 167/300
successes versus 169/300 with full Bo4 and 144/300 with one proposal.
Four-way calls fell from 3,099 to 1,500 (51.6\%) and generated chunks
from 12,396 to 7,608 (38.6\%). First-ten agreed with full Bo4 on all
300 terminal outcomes. First-seven versus random-seven was
$+3.67$ pp $[-2.00,9.33]$. Those policies realized 2,081 and 1,567
calls, respectively, leaving timing effects at equal realized cost
unresolved.

On 200 new Object layouts disjoint from both the pilot and calibration
cohorts, first-five and full Bo4 each succeeded 110 times
($0.00$ pp $[-2.00,2.00]$), with two outcome changes in each direction.
Calls fell from 2,061 to 1,000. Relative to the single-proposal control
(98/200), the difference was $+6.00$ pp $[-2.00,14.00]$.

The delayed control used one proposal at queries 1--5 and Bo4 thereafter. On calibration, Late succeeded 141/300 times and used 1,757 four-way calls, versus 167/300 and 1,500 calls for first-five. Its change from single-proposal control was -1.00 pp [-2.67, 0.67]. First-five versus Late was +8.67 pp [2.33, 15.00]. On fresh starts, Late succeeded 98/200 times and used 1,136 calls, versus 110/200 and 1,000 for first-five. First-five versus Late was +6.00 pp [-2.00, 14.00]. Late showed no positive point-estimate gain over one proposal on either cohort.

Across 500 paired Object layouts, first-five outperformed Late (+7.60 pp [2.60, 12.60]), with 2,500 versus 2,893 realized four-way calls. The combined percentile interval resamples physical layouts within task and cohort. Calibration helped choose the cutoff, and the fresh comparison is reported separately. Realized calls depend on episode termination. We did not match realized call counts exactly.

We evaluated the proportional rule prospectively on ten LIBERO-Goal tasks, using 300 official starts paired across all four arms. With a 300-step cap and at most 20 actions per query, the preregistered cutoff was $K=\operatorname{round}(15/3)=5$. We used the released Goal critic and the shared actor. The four arms (single, full Bo4, EWS-5 and Late-5) succeeded 145, 165, 159, 138/300 times, respectively. EWS-5 minus full was -2.00 pp [-6.00, 2.00], with 1,484 versus 2,996 four-way calls. Late minus single was -2.33 pp [-6.33, 1.67], with 1,751 four-way calls. The one-third rule met both prespecified point-estimate criteria on Goal. Together with the separately set Object and Long Task2 windows, this supports early windows across the three evaluated suite settings. The result does not establish a universal cutoff. Initial physics and seeds were identical across arms, while the first wrist image could differ by up to two pixel levels. EWS and full Bo4 chose the same first winner on 293/300 starts. Goal provides a paired schedule comparison, with rendering variability in the common prefixes.

For Long Task2 we set a separate eight-query window from development
observations of the longer task horizon. The two cutoffs use roughly a third of their query caps (5/14 and 8/26). First-eight and full Bo4
succeeded in 127/256 and 124/256 episodes, respectively
($+1.17$ pp $[-5.08,7.42]$). Relative to the single-proposal control
(93/256), the difference was $+13.28$ pp $[7.03,19.53]$. Calls fell
from 4,949 to 2,048 (58.6\%) and chunks from 19,796 to 11,077
(44.0\%). Both confirmation comparisons are within the 3 pp
point-estimate tolerance. The schedule reduces the number of
four-way queries, and prefix sharing reduces work within a query.

\input{\paperdir tables/episode-cost.tex}
Table~\ref{tab:episode-cost} sums the recorded core decision times within each episode. On the Object calibration cohort, sharing plus EWS used 2,201 ms of mean inference per episode versus 2,282 ms for one proposal ($-3.5$\%), with 167 versus 144 successes. On fresh Object starts the corresponding means were 2,184 and 2,260 ms with 110 and 98 successes. These totals also reflect changes in the number of queries reached. On Goal, EWS-5 used 2,943 ms of mean core inference per episode, 13.4\% above the single-proposal arm (2,595 ms). The Object result on core time did not transfer to Goal.

\subsection{Decision latency and critic storage}
Across Object rollout decisions, mean latency is \BaselineTime{} ms for one
proposal, \OriginalTime{} ms for original Bo4 and \SharedTime{} ms for shared
Bo4. Table~\ref{tab:scaling} uses a separate fixed-observation protocol with
the critic active at every candidate count. Table~\ref{tab:precision} reports
winner agreement and paired Object outcomes with packed critic matrices.
The storage counts include scales. Computation uses unpacked weights.

\input{\paperdir tables/graph.tex}

\input{\paperdir tables/profile.tex}

\input{\paperdir tables/precision.tex}

\subsection{Denoising schedule}
Figure~\ref{fig:stages} and Table~\ref{tab:profile} show that action generation
dominates eager stage time, while vision and prefix account for a substantial
fraction of registered FLOPs. Logical activation and parameter bytes provide
modeling inputs. Physical DRAM traffic remains unmeasured.

Graph replay reduces expert time by \ExpertReduction\% at four candidates.
Figure~\ref{fig:decision} and Table~\ref{tab:graph} show that the latency
increase with candidate count is more visible after replay. The reduction
captures graph-removable elapsed time on this software stack, including host
dispatch and device scheduling effects. In the \FpChecks{} strict FP32
comparisons, actions, scores and selected indices agreed.

\begin{figure}[tb]
\centering
\includegraphics[width=\columnwidth]{\paperdir figures/decision-latency.pdf}
\caption{Matched eager and CUDA Graph decisions on a fixed Object observation.
Input refresh is included. One-time capture is excluded.}
\label{fig:decision}
\end{figure}
\begin{figure}[tb]
\centering
\includegraphics[width=\columnwidth]{\paperdir figures/stage-latency.pdf}
\caption{Eager stage profile on A100. Stacks sum separately measured stage medians.}
\label{fig:stages}
\end{figure}

\section{Discussion and Limitations}
The flat eager curve describes the measured execution regime. Replay removes
substantial scheduling cost and reveals increased incremental cost for wider
batches. Determining how much comes from individual launches would require
kernel-level traces. All timing is on A100. Target-chip latency, bandwidth
and energy remain open. Closed-loop success measurements use eager execution.
Three observations are consistent with a shorter opportunity to recover
late in an episode. In a prior Long Task2 timing study, only one episode
that reached query 16 ultimately succeeded. On Object, Late used more
four-way calls than EWS-5 and yielded no positive gain over one proposal
by point estimate. Candidate zero's median twin-min score rose from
$-56.41$ at query 1 to $-32.37$ at query 6 on 3,099 archived full-Bo4
queries (Fig.~\ref{fig:q0-gate}). A pooled lower-quartile gate at
$-53.78$ would mark 747 of its 775 queries in the first five. Late
samples condition on reaching that query. Compounding early errors can
also make late trajectories harder to rescue. The remaining-horizon
interpretation is consistent with the reach curve and the delayed
control, and these observations do not isolate it from compounding
error. Both mechanisms motivate testing fixed early windows. The score
analysis is retrospective, and earlier visual-residual and
critic-advantage gates showed no reliable incremental benefit in their
tested settings.
Exploratory score-gated windows and disagreement-penalized selection
yielded lower success point estimates than the fixed window in their
respective matched simulations.

\begin{figure}[tb]
\centering
\includegraphics[width=\columnwidth]{\paperdir figures/q0-gate.pdf}
\caption{Candidate-zero score drift and retrospective low-score gate
rates on shared full-Bo4 Object trajectories. Percentile thresholds
are pooled over all 3,099 recorded queries. Gate rates condition on
reaching each query index.}
\label{fig:q0-gate}
\end{figure}

VGAS supplies the actor, critic and selection rule. Object and Long use the
same LIBERO simulator \cite{libero}. Long is scoped to Task2 selected in
development. Object and Long use separately calibrated windows. Paired
intervals are descriptive. On Long Task2, the single-proposal control
succeeded in 38/128 and 55/128 episodes under the two noise banks. The 3 pp point-estimate rule guides schedule
selection, while formal equivalence requires a prespecified margin and test.
The measured inference totals omit environment execution and action postprocessing. Complete episode wall time and energy are unmeasured. In an offline Object probe of 150 four-proposal
queries, step-four estimates matched the winner from the complete ten-step
denoising run in 52\% of cases, below the 80\% screening threshold. The
reported workload therefore retains all ten denoising steps per proposal.

\section{Conclusion}
Shared observation encoding reduces A100 decision time for frozen VGAS
selection. We observed that using four proposals only at early queries
retained the full schedule's observed success within the specified
point-estimate tolerance on new Object layouts and on Long Task2. Prefix
reuse and early-window selection address work within a query and across
queries, respectively. Measuring complete episode wall time would further characterize the deployment impact.

\begingroup\small
\input{\paperdir references.tex}
\endgroup
\end{document}
